\section{Introduction}
\label{sec:introduction}

\subsection{Abstract}
\label{subsec:abstract}
\textbf{Background:} Designing responsive, cross-device IoT games on resource-constrained microcontrollers requires balancing low-latency local execution with secure, asynchronous network communication.
\textbf{Objective:} This project presents BiteBound, a physical, tilt-controlled labyrinth game built on a Waveshare ESP32-S3-Touch-LCD-1.69 featuring real-time communication with two remote dashboards.
\textbf{Methods:} The system utilizes a dual-core ESP32-S3 microcontroller to handle network communication, sensor inputs, physics simulation, game logic and display rendering, while a Kotlin-based Android app and a Node-RED interface provide bidirectional control via telemetry and command messages over the Message Queuing Telemetry Transport (MQTT) protocol.
\textbf{Results:} The firmware, running on a network and game logic core, achieves a stable 50\,Hz frame rate. A bakery-themed frontend provides a playful user experience, while the companion interfaces successfully allow for remote monitoring and control of the game state during gameplay.
\textbf{Conclusion:} BiteBound demonstrates a robust, responsive end-to-end architecture bridging physical embedded interactions with digital companion interfaces.

\subsection{Mission Statement}
\label{subsec:mission_statement}
The objective of the BiteBound project is to implement a highly responsive, low-latency gaming platform demonstrating efficient multi-core task scheduling on resource-constrained hardware. Centered on a playful, bakery-themed aesthetic, the system requires navigating a virtual cherry (ball) through procedural mazes. The main technical challenge is maintaining a stable, high-frame-rate rendering pipeline alongside bi-directional MQTT communication, supported by native Android and Node-RED dashboards to deliver a unified, multi-platform user experience.

\subsection{Document Structure}
\label{subsec:document_structure}
The remainder of this report is organized as follows:
\Cref{sec:technical_concept} outlines the high-level system architecture.
\Cref{sec:hardware} details the firmware, sensor integration, physics, game logic, and display optimizations.
The companion interfaces are presented in \Cref{sec:android_app} (Android App) and \Cref{sec:nodered} (Node-RED).
Finally, \Cref{sec:project_management} covers project organization and milestones, while \Cref{sec:discussion, sec:summary-future-work} provide evaluation, conclusions, and future outlook.